\section*{ВВЕДЕНИЕ}
\addcontentsline{toc}{section}{ВВЕДЕНИЕ}

\textbf{Актуальность темы.} В современных условиях цифровизации бизнеса и повсеместного перехода на удаленный или гибридный формат работы, эффективное управление проектами становится критически важным фактором успеха любой организации. Использование традиционных методов учета задач (электронные таблицы, мессенджеры) приводит к потере информации, срыву сроков и снижению производительности команды. Разработка специализированных веб-приложений для трекинга задач позволяет централизовать рабочие процессы, автоматизировать контроль поручений и обеспечить прозрачность работы на всех этапах жизненного цикла проекта.

\textbf{Степень разработанности проблемы.} На рынке существует множество систем управления проектами (Jira, Trello, Asana и др.). Однако большинство корпоративных решений обладают избыточным функционалом, сложны в освоении и имеют высокую стоимость лицензий, что делает их нерентабельными для малого бизнеса, студенческих команд или индивидуального использования. В связи с этим, создание легковесного, интуитивно понятного веб-приложения с базовым, но достаточным набором функций (CRUD-операции, ролевая модель, прикрепление файлов) остается актуальной практической задачей.

\textbf{Объект исследования} — процессы управления проектами, распределения задач и контроля их выполнения в рабочих группах.

\textbf{Предмет исследования} — программные средства, алгоритмы и методы веб-разработки, применяемые для создания систем отслеживания задач.

\textbf{Целью данной работы} является проектирование и реализация клиент-серверного веб-приложения, предназначенного для эффективного управления проектами и отслеживания статусов выполнения задач с учетом ролевой модели доступа.